\chapter{总结与下一步工作}

这轮优化最重要的成果并不是某一个百分比，而是建立了一条可以继续使用的分析路径。我们先用完整标记确认 workload 真实完成，再用 pc-hot 看指令分布、wait-hot 看核心空闲与阻塞、syscall profile 看 ABI 调用形态；得到机制假设后只改变一个主要因素，最后回到固定条件的完整 A/B。这个过程使一些局部漂亮但端到端无效的方案能够及时回退。

最终保留下来的改动也有清晰的共同点。syscall 和 VFS 复用一次调用中已经取得的参数、fd 与 lookup 结果；文件系统缓存根据实际复用决定准入和容量；ELF 与只读 mmap 在不同进程间共享稳定物理页；mprotect、mmap、munmap 和 COW 只对真实页表变化执行必要同步；匿名页和页表页的清零则由 idle CPU 提前完成。它们没有改变用户可见 ABI，而是减少了同一语义的重复实现成本。

现有结果仍有边界。历史实验跨越多个 main 基线，不能把独立百分比相加；部分结果只有单轮完整记录；LoongArch 与 RISC-V 的 Linux 对照条件尚未完全统一。因此，最终提交前仍应在固定宿主 CPU、相同 QEMU 和镜像上重跑当前 main，给出 WaterOS 与 Linux 的统一基线、多轮中位数和离散程度。

下一步优先关注三项工作。第一，重新 profile 当前主线，确认只读页共享和零页池之后的新热点，避免继续优化已经消失的路径。第二，检查 rustc 关键阶段的单核瓶颈与 vCPU 不均衡，区分 scheduler、锁等待和 guest 本身串行阶段。第三，在正确性可控的前提下研究 VirtIO 多请求并行与更精确的 TLB range shootdown；历史分支中相关方案已有失败记录，新的实现必须先解释中断、完成队列和错误恢复边界。

性能优化不会在“找到最热函数”时结束。真正可维护的结果应当同时满足三个条件：能够从 workload 数据解释根因，能够在接口层说明正确性边界，并能够在完整运行中重复观察到收益。本文记录的工具、反例和复现方法，就是为了让后续工作仍然遵守这三个条件。
